\documentclass{article}
\usepackage[utf8]{inputenc}
\usepackage[T1]{fontenc}
\usepackage{amsmath,amssymb,amsthm}
\usepackage{hyperref}
\title{Finding report: Polynomial calculus degree at least 5 over F2 for the onto functional pigeonhole principle with surplus 8}
\author{Lindheim, Björn}
\date{}
\begin{document}
\maketitle
\begin{center}DOI: 10.5281/zenodo.24056805\end{center}
\section{Introduction}
\label{sec:pc-degree-5-onto-fphp-surplus-8}

Cite as: https://doi.org/10.5281/zenodo.23996615 (all versions; this version: https://doi.org/10.5281/zenodo.24056805)

\subsection*{Setting}

The onto functional pigeonhole principle $\mathrm{onto\text{-}FPHP}^{n+k}_n$ says that $n+k$
pigeons can be mapped one-to-one onto $n$ holes; it is false for every $k\ge1$. In Boolean
variables $x_{ij}$ (pigeon $i$ in hole $j$) its axioms are $\sum_j x_{ij}-1$,
$\sum_i x_{ij}-1$, $x_{ij}x_{i'j}$ and $x_{ij}x_{ij'}$, the linear encoding used for this family
\cite{beame1998counting,razborov2001php}. These results are for the linear encoding only: translating the CNF clause by clause instead turns each pigeon clause into a polynomial of degree $n$ \cite{razborov2001php}, so degree bounds below $n$ say nothing about that encoding. A polynomial calculus (PC) refutation over
$\mathbb F_2$ derives $1$ from the axioms and the Boolean axioms $x^2-x$ by linear combinations
and multiplication by a variable; its degree is the largest degree of a line.

Over $\mathbb F_2$, with surplus $k=2^\ell$, the PC degree is exactly 2 at $k=2$ and exactly 4 at
$k=4$ (companion results in the same library). This finding is about $k=8$.

\subsection*{Result}

For every $n\ge7$, $\mathrm{onto\text{-}FPHP}^{n+8}_n$ has no PC refutation over $\mathbb F_2$ of
degree 4. Its PC degree is therefore at least 5, while the general bound for surplus divisible
by $p^2$ (companion finding) gives only 3 at $p=2$. So at $p=2$ the third level is strictly
harder than the second for every $n\ge7$.

The proof combines a general reduction with finite certificates: small ones for levels 0--2 and,
for level 3, one of 95,550 steps that the Lean kernel checks in 2,154 pieces, plus 60 coverage checks.

\subsection*{Why it matters}

Razborov's Question 7 asks whether PC degree for surplus $p^\ell$ stays bounded in characteristic
$p$ \cite{razborov2001php}. Riis and Sitharam's Theorem 33 gives degree at least 7 for $k=8$ once
$N\ge196$ \cite{riis1997nonconstant}, which is stronger than this result for large $n$. The range
$7\le n<196$ is not covered by that theorem. This result covers it, and it is machine-checked.

\subsection*{Proved vs checked}

The result is proved in Lean 4 for every $n\ge7$; the certificate is part of the proof, not a
separate check up to a bound. The proof uses only the standard axioms \texttt{propext},
\texttt{Classical.choice} and \texttt{Quot.sound}. A second agent in the run checked the certificate
independently, with its own program written from the Lean definitions, not from the author's
code. That is an independent check of the certificate, not an independent Lean re-proof. No
person has reviewed the proof. The same library also states the result for standard (syntactic)
PC (\texttt{std\_no\_deg4\_k8}).

The Lean theorem excludes refutations of degree exactly 4; smaller degrees are excluded by monotonicity in
the degree bound, stated in the library as \texttt{Programme.«maintenance-20261108».DegreeMono.no\_deg\_le4} (and \texttt{deg\_ge5}).

\subsection*{Novelty statement}

This result extends:

\begin{itemize}
\item Riis-Sitharam, ECCC TR97-048, Theorem 33: PC degree \textgreater{}= 2\textasciicircum{}l - 1 for onto-PHP\textasciicircum{}\{N+p\textasciicircum{}l\}\_N over characteristic p once N \textgreater{}= (21+3p\textasciicircum{}2)(2\textasciicircum{}l - 2) - p (same family; proof not checked by us; at l = 2 this is \textgreater{}= 3 against our \textgreater{}= p+1 from n = 2p-1; at p = 2, l = 3 it is \textgreater{}= 7 for N \textgreater{}= 196 against our \textgreater{}= 5 for n \textgreater{}= 7).
\item Beame-Riis 1998, Theorem 8.1 and Lemma 8.7: Nullstellensatz degree between 2\textasciicircum{}l - 1 and p\textasciicircum{}l - 1 (multiplier degree) for the same family; Section 10 lists polynomial-calculus lower bounds as open.
\item Riis 1993, as Theorem 8 of Razborov's survey: degree-d PC refutations when C(m,d) and C(n,d) differ mod p, which gives the matching upper bound p for k = pc, p not dividing c (via Lucas).
\item Razborov, Proof complexity of pigeonhole principles (2001), Question 7: the fixed-p, fixed-l bounded-degree question for onto-FPHP\textasciicircum{}\{n+p\textasciicircum{}l\}\_n.
\item Miksa-Nordstrom (JACM 2024): general PC degree method; onto-FPHP with m = n+1 easy (Riis); the case of characteristic dividing the surplus not covered.
\item BGIP 2001: proves linear PC degree for Tseitin mod p over characteristic q containing a primitive root of unity of order p (Thm 15) and for MOD\_p\textasciicircum{}n over F\_q, q not dividing p (Cor 20: ``requires degree \textgreater{} delta n''). This is the cross-characteristic regime. Our results concern the same-characteristic regime p $\mid$ k, so there is no overlap. The paper does not treat onto-PHP.
\item Alekhnovich-Razborov 2003: ``immunity'' method giving linear PC degree bounds for mod\_p Tseitin over characteristic different from p, for flow tautologies and for random k-CNF in characteristic 2 (abstract only). No PHP or onto-PHP result. The method is relevant only as possible technique.
\item Impagliazzo-Pudlak-Sgall 1999: simpler proof of Razborov's ceil(n/2)+1 PC lower bound for (injective, non-onto) PHP\textasciicircum{}m\_n over any field, plus an upper bound in large characteristic (abstract). This is a different principle: without the onto and functional axioms the bound is independent of p. It does not transfer to the exact degree p or p+1 for onto-FPHP.
\item BIKPRS 1996/97: Nullstellensatz lower bounds for Count\_q over Z\_m with q not dividing m (Thm 3.2) and over Q (Thm 4.6), plus the PC-to-NS simulation (Thm 5.4). This is Nullstellensatz and Count, not PC degree for onto-FPHP, and the paper excludes the case q $\mid$ m.
\item Grigoriev FOCS 1998 (snippet only, not opened): Nullstellensatz degree for Tseitin mod p over other characteristics. It is a predecessor of BGIP; NS, not PC, and a different principle.
\item Filmus-Lauria-Mikša-Nordström-Vinyals 2025: PC space and separations. It mentions the functional PHP only for space limits, not for characteristic-dependent degree.
\item Ken 2022/2024: restates Razborov's injPHP bound ``no PC-refutation ... with degree \textless{}= m/2'' (Thm 24) and extends it to NS over finite rings (Cor 25). Onto-PHP appears only in reductions (Props 12, 15), and no degree bound for it is proved.
\item Reduction transfer: none of the opened sources gives a low-degree reduction from Count\_q, Tseitin mod q or Mod\_q (q != p) to onto-FPHP\textasciicircum{}\{n+k\}\_n with p $\mid$ k over F\_p. Since the cross-characteristic lower bounds need q != p, they do not yield results 1-3 directly.
\end{itemize}

Our prior-art search (2026-11-06) covered:

\begin{itemize}
\item Riis, Sitharam, Non-constant degree lower bounds imply linear degree lower bounds, ECCC TR97-048 (1997), pp. 36-38 and 45-46 (Theorem 33, Propositions 27 and 28), opened 2026-10-29
\item Beame, Riis, More on the relative strength of counting principles, DIMACS Series 39 (1998), homes.cs.washington.edu/\textasciitilde{}beame/papers/riis.pdf: Theorem 8.1, Lemma 8.7, Section 10, opened 2026-10-27 and 2026-11-01
\item Razborov, Proof complexity of pigeonhole principles, DLT 2001 (LNCS 2295), people.cs.uchicago.edu/\textasciitilde{}razborov/files/php\_survey.pdf: Theorem 8, Question 7, footnote 5, opened 2026-10-27 and 2026-10-29
\item Riis, BRICS Report RS-99-29, opened 2026-10-29
\item Miksa, Nordstrom, A generalized method for proving polynomial calculus degree lower bounds (CCC 2015; JACM 2024; arXiv 1505.01358), opened 2026-10-27
\item de Rezende et al., Theory of Computing 21(6) (2025), opened 2026-10-29
\item Lauria et al., Theory of Computing 21(4) (2025), opened 2026-10-27
\item Buss, Grigoriev, Impagliazzo, Pitassi, ``Linear gaps between degrees for the polynomial calculus modulo distinct primes'', JCSS 62(2):267-289, 2001 (STOC 1999). Opened: journal final PDF, https://mathweb.ucsd.edu/\textasciitilde{}sbuss/ResearchWeb/lineargaps/journalfinal.pdf
\item Alekhnovich, Razborov, ``Lower bounds for polynomial calculus: non-binomial case'', Proc. Steklov Inst. Math. 242 (2003) 18-35 (FOCS 2001). Opened: abstract only, https://mathnet.ru/eng/tm403 (full text not opened).
\item Impagliazzo, Pudlak, Sgall, ``Lower bounds for the polynomial calculus and the Groebner basis algorithm'', Computational Complexity 8(2):127-144, 1999 (ECCC TR97-042). Opened: ECCC abstract only, https://eccc.weizmann.ac.il/report/1997/042/
\item Buss, Impagliazzo, Krajicek, Pudlak, Razborov, Sgall, ``Proof complexity in algebraic systems and bounded depth Frege systems with modular counting'', Computational Complexity 6(3):256-298, 1996/97. Opened: full PDF, https://mathweb.ucsd.edu/\textasciitilde{}sbuss/ResearchWeb/nullstellensatz/token.pdf
\item Grigoriev, ``Tseitin's tautologies and lower bounds for Nullstellensatz proofs'', FOCS 1998, pp. 648-652. Snippet only, not opened (dblp/researchr entries).
\item Mikša, Nordström, ``A generalized method for proving polynomial calculus degree lower bounds'', CCC 2015 (arXiv 1505.01358). Already credited. Seen in search results; not re-opened here.
\item Filmus, Lauria, Mikša, Nordström, Vinyals, ``Towards an understanding of polynomial calculus: new separations and lower bounds'', Theory of Computing 21(4), 2025 (ICALP 2013). Opened: abstract page, https://zesty2.cs.uchicago.edu/articles/v021a004
\item Eitetsu Ken, ``On some Sigma\_0\textasciicircum{}B-formulae generalizing counting principles over V\textasciicircum{}0'', arXiv 2203.10237 (2022; a journal version appears to be Arch. Math. Logic 2024, https://link.springer.com/article/10.1007/s00153-024-00938-1). Opened: ar5iv HTML and Springer page, first \textasciitilde{}100k characters of each.
\item Web searches for 2000-2026 papers on PC degree of onto-PHP or Count where the characteristic divides the surplus: no further hits beyond the items above.
\end{itemize}

Prior-art answer recorded by the owner (basis stated in the file); published in \texttt{prior-art/}.

If you know of related work we missed, please open an issue.

This result was found by AI research agents (the autoresearch harness). It was proved in Lean; the verification steps below build the published programme library.

\subsection*{How it was produced}

The declaring Lean sources are in \texttt{checks/library/}.

\subsection*{How to verify}

Lean toolchain: \texttt{leanprover/lean4:v4.35.0}. \url{https://github.com/kodakodapa/autoresearch-findings/tree/ad1259e5894f224a3d1eda62a1914da4b4dd699e/programmes/proof-complexity}. Install elan and open a new shell after installing elan, then pin the library subtree while keeping the finding checks:

\begin{verbatim}
git clone https://github.com/kodakodapa/autoresearch-findings.git
cd autoresearch-findings
git checkout ad1259e5894f224a3d1eda62a1914da4b4dd699e -- programmes/proof-complexity
cd programmes/proof-complexity && lake build 'Programme.«20261028-0841-proof-complexity».L3.Main'
AR_LIBRARY=$PWD python -I ../../claims/pc-degree-5-onto-fphp-surplus-8/checks/lean-K18.py
\end{verbatim}

A full \texttt{lake build} builds the whole library. The named targets build only the claim modules and their dependencies. Indicative measurements from the 2026-11-06 reproduction on 2 cores / 7.8 GB: K13--K15 90 s; K1--K6 440 s; K18 87 min (90 certificate modules). Peak memory was 4.7 GB for one Lean process. Build times depend on hardware and cached dependencies.

\texttt{checks/run-copies/}: the author's working copies from the run; they import the run's private \texttt{Work} package and are not built standalone.

\section{Claims}

\subsection*{Claim K18}

Lean \texttt{Work.L3.Main.no\_deg4}: for every $n\ge7$, \texttt{ontoFPHP n 8} (onto-FPHP$^{n+8}_n$, functional encoding) has no degree-4 \texttt{PPCRefutation 2} (polynomial calculus over $\mathbb Z/2$, line degree 4). Unconditional.

Status: proved (axioms: propext, Classical.choice, Quot.sound).

Re-checked at publication on 2026-11-08: pass (\texttt{results/publication/lean-K18.json}).

\paragraph{Verified.}
No check scripts.

\paragraph{Evidence.}
\begin{itemize}
\item library check \texttt{checks/lean-K18.py} (library theorem \texttt{Programme.«20261028-0841-proof-complexity».L3.Main.no\_deg4})
\item proof \texttt{checks/run-copies/K18-Main.lean}: Work.L3.Main.no\textbackslash{}\_deg4 --- pass (\texttt{results/K18-4.json})
\end{itemize}

\paragraph{Prior art.} Partial prior art.


\paragraph{Provenance.}
Run \texttt{20261028-0841}, claims K18, run-folder commit \texttt{19b6167239a90c29b4a3f0ff0f64c79016faf2f8}, published 2026-11-06.

\section{Discussion}

The upper bound from Riis's criterion is 8 \cite{riis1993independence,razborov2001php}, so for
$k=8$ the degree lies between 5 and 8 for $n\ge7$. This is one case, not a growth result in
$\ell$. We have not checked the proof of Theorem 33 and cite it for comparison only.

\section{How to verify}

For programme findings, follow the clone, build and run steps in each finding. The declaring Lean sources are in \texttt{checks/library/}; the publication checks import their modules and print each theorem's axioms. The output of the checks re-run at publication is in \texttt{claims/<slug>/results/publication/}. The checks are Python scripts and need the packages of the research image (the image's \texttt{requirements.txt}); the Lean checks also need \texttt{lean}.

\section{Revisions}

\begin{itemize}
\item 2026-11-06 --- Math rendering on the site; paper lists only cited references
\item 2026-11-06 --- Library published buildable; verification steps; run copies labelled
\item 2026-11-06 --- Verification steps; stale run copies removed; encoding stated
\item 2026-11-06 --- Publication re-checks and reproofs published; prior art extended; texts corrected
\item 2026-11-06 --- Library theorem names and checked statements shown; reproof wording corrected; build cost
\item 2026-11-08 --- Certificate wording; degree-monotonicity theorems named (degree \textgreater{}= 5)
\end{itemize}

\bibliographystyle{plain}
\bibliography{refs}

\end{document}
